\chapter{Results and Evaluation}
\label{ch:results}

\section{Working Prototypes}
\label{sec:results_prototypes}

\subsection{Drums Real-Time System}
\label{sec:results_drums}
% TODO: drums_v1 ep755 in Pure Data. Subjective: responds clearly to voice
% and beatbox input. Quantitative metrics (paste from session diagnostics):
%   Encoder differentiation: latent diff = 2.40 for random-noise inputs.
%   Decoder responsiveness: forward diff = 0.0045 between distinct inputs.
%   Silence handling: silence amplitude 0.004 < signal amplitude 0.003 only
%     marginally; the model produces drum-like output regardless of input
%     but varies recognisably with the input's transient content.
% Latency: sub-10 ms measured (see Section \ref{sec:impl_realtime}).
% Note: drums_v2 (with official v2 architecture) trained later — append
% its metrics here.

\subsection{Guitar Autoencoder (in-distribution)}
\label{sec:results_guitar_ae}
% TODO: guitar_v6_best (epoch 1935, version_4 of guitar_v6_bbcc6d36c3).
% When fed real guitar audio from GuitarSet:
%   Output amplitude ratio: 0.36-1.17 across samples.
%   STFT magnitude correlation: 0.5+ on at least one test sample.
%   Subjective listening: recognisable as guitar; pluck-like attacks and
%     identifiable note content; degraded compared to input.
% This is the first model in the series that demonstrably works as a
% guitar autoencoder. Include reconstruction WAV examples in Appendix B.

\section{Failure Taxonomy across Six Guitar Iterations}
\label{sec:results_failures}

% TODO: For each version, a sub-subsection with:
%   - Configuration (dataset, latent size, special changes)
%   - Failure observation (perceptual + metric)
%   - Diagnostic test (TensorBoard graphs, Python diagnostic outputs)
%   - Lesson carried forward
% A summary table at the end consolidates the taxonomy.

\subsection{guitar\_v1: Posterior Collapse}
% Beta=0.1 constant from step 0 -> KL decreasing -> encoder ignored.
% Diagnosis via TensorBoard regularisation curve (drops 0.5 -> 0.25).
% Python test: encoder diff = 1.25 between silence and 440 Hz sine,
% but decoder output diff = 0.006 (decoder ignores latent).
% Fix carried forward: log-space beta warmup.

\subsection{guitar\_v2: Partial Collapse from Oversized Latent}
% LATENT_SIZE = 128, beta warmup correct. Total KL plateaued at ~0.6 i.e.
% 0.005 nats/dim. Encoder using almost none of the capacity.
% Diagnosis: Python test forward diff = 0.003 -> decoder ignores latent
% variation. Output amplitude only 0.012 = 60x quieter than input.
% Fix carried forward: LATENT_SIZE = 16.

\subsection{guitar\_v3: Partial GAN Collapse}
% LATENT_SIZE = 16. KL rose to 0.65 (0.04 nats/dim - much better).
% Encoder responding (forward diff = 0.042 at peak). But Phase 2 GAN
% pushed amplitude down and never recovered. Output was quiet
% noise that responded slightly to input but did not sound like guitar.

\subsection{guitar\_v4: Discriminator Dominance}
% Cleaner mic-only 3.05h dataset. KL = 0.80. But Phase 2 collapsed
% completely:
%   pred_real = 3.34, pred_fake = -2.31, gap = 5.65.
%   Output amplitude (0.0045) < silence amplitude (0.017).
% The decoder gave up; it produced more sound for silence than for
% signal -- a textbook adversarial mode collapse.
% Counter-intuitive lesson: cleaner data made the failure WORSE because
% the discriminator could memorise more easily.

\subsection{guitar\_v5: Discriminator Over-Weakening}
% Reduced discriminator capacity 64 -> 32, layers 4 -> 3. Hoped to
% prevent v4-style dominance. Instead the discriminator collapsed:
%   pred_real = -0.90, pred_fake = -1.70, gap = 0.80.
% Discriminator labels real audio as fake. Generator has no useful
% adversarial gradient.

\subsection{guitar\_v6: Official RAVE v2 + Causal}
% Abandoned custom c16_r10. Used official v2.gin + causal.gin with
% two minimal overrides (LATENT_SIZE = 16, PHASE_1_DURATION = 1.5M).
% Phase 1 produced normalised latent (mean = 0.11, std = 0.82). Phase 2
% gap peaked at 11.88 then DECREASED to 8.55 then plateaued around
% 5-9 -- the FIRST training in our series where the GAN gap reversed
% direction. The model is a working guitar autoencoder.
% Voice input: still produces constant noise -- but for a different
% reason (decoder mode collapse on OOD input, not GAN failure).
% This is the critical scientific distinction: the architecture works,
% the input domain mismatch does not.

\subsection{Summary Table}
% TODO: 6-row table with columns:
% Version | Dataset | Latent | Outcome | Key Metric | Lesson

\section{Voice Input as Out-of-Distribution}
\label{sec:results_voice_ood}

% TODO: This is the most important section for explaining why
% voice-to-guitar failed even with a working guitar autoencoder.
%
% Quantitative test on guitar_v6:
%   - Input "voice 220 Hz with formants": encoder latent z mean = -550,
%     std = 2484 (vastly out of distribution from training data where
%     z mean is near 0 and std near 1).
%   - The encoder extrapolates outside the trained manifold.
%   - The decoder produces near-zero or near-constant output for these
%     extrapolated latents.
%
% Interpretation: BRAVE is trained to reconstruct guitar. Voice, while
% it has pitched content, has a fundamentally different spectral
% envelope (formants, vocal tract resonances). The encoder maps voice
% to latent points outside the guitar manifold, where the decoder is
% not trained to produce meaningful output.
%
% This is a known limitation in the timbre-transfer literature; we
% confirm it quantitatively for our specific BRAVE + guitar setup.

\section{Real-Time Performance}
\label{sec:results_latency}

% TODO: Measured latency in Pure Data. Block size 64 at 44.1 kHz = 1.45 ms.
% Plus driver buffer ~5 ms (WASAPI). Plus model forward < 1 ms. Total ~7 ms.
% Confirms the BRAVE design target of sub-10 ms.
